\section{Discussion}
\label{sec:discussion}

\subsection{What the ablation identifies}
The four hybrid conditions were designed so that their pairwise contrasts each
answer one question, and on this data every answer is the same.
\textsc{Trained} against \textsc{Frozen} asks whether optimising the circuit
matters: the per-cell Sharpe ratios correlate at \num{1.000}, so it does not.
\textsc{Frozen} against \textsc{Surrogate} asks whether the particular random
feature a quantum circuit computes differs from a classical one of equal
parameter count and matching output scale: on no panel can the two be told
apart, and the sign of the difference changes between panels. \textsc{Surrogate}
against \textsc{None} asks whether any extra scalar feature helps the head at
all: it does not. The circuit is a random feature, the random feature is not
special, and the head does not need it.

The same three contrasts give the same three answers on the US panel
(Section~\ref{sec:results:us}), where gross performance is negative rather than
positive. A result that survives a sign change in the underlying market is
about the layer, not about the window.

A null result on a task where nothing works might say little about the layer.
The locked holdout is the positive control for that concern. On those \num{135} days the shared signal is profitable for all
nine conditions, with gross annualised Sharpe ratios from \num{+0.42} (MLP)
to \num{+1.83} (the no-circuit control), and the ablation is unchanged there: trained beside frozen, the
surrogate no worse than either, and no contrast significant. The circuit adds
nothing in the regime where the features fail and nothing in the regime where
they work.

The result does not show that quantum convolutional layers are useless. It
concerns this task: daily cross-sectional ranking from twenty days of
technical features, where the achievable information coefficient is around
\num{0.006} and the spread of one model's Sharpe ratio across seeds
(standard deviation \numrange{0.13}{0.61}) is as large as the spread between
models (\num{0.17}). A layer cannot demonstrate value in a regime where the
signal it would need to improve on is itself indistinguishable from noise.

\subsection{Why earlier numbers looked different}
The reference study this work audits \citep{chen2025quantumb} reports an
out-of-sample Sharpe ratio of \num{0.538} $\pm$ \num{0.042} for its quantum
model, about \SI{72}{\percent} above its best classical baseline, and defines
that ratio as annualised. Its accompanying public implementation, which is the
artefact we re-ran, computes the ratio as mean over standard deviation with no
$\sqrt{252}$ factor and yields \num{0.33} on the same window; on the
conventional scale that is \num{5.2}. Whichever figure one takes, four things,
each individually sufficient, separate it from Table~\ref{tab:headline}. It
was computed on \num{22} trading days, a window over which the bootstrap
interval spans several Sharpe units in either direction; the authors
themselves note that \num{22} days give limited power and that about \num{45}
would be needed for \SI{90}{\percent} power, and defer a full-sample backtest
to future work. The headline is a single-seed model selection (noise
experiments use five runs), and the seed draw alone moves a fold's Sharpe by
up to \num{2.2} units for the two conditions in Figure~\ref{fig:seeds} and by
\num{4.39} for LightGBM. The split adjustment
ran forward in time, corrupting every momentum and volatility feature with
artefactual returns of up to \num{2300}\%. And the paper's own encoder-only
ablation already scores \num{0.412} $\pm$ \num{0.039}, so most of the reported
edge over the Transformer baseline is the temporal encoder, not the circuit, the same decomposition our \textsc{None} condition makes explicit. None of
these is exotic; all are common in the applied quantum machine learning
literature \citep{shen2026quantum,bowles2024better}, where evaluation practice
has not yet caught up with the standards of empirical finance.

\subsection{The circuit is trainable; that is not the problem}
The diagnostics rule out the explanation a quantum-computing reader would reach
for first. Gradient variance is flat from two to twelve qubits, so there is no
barren plateau; expressibility improves with depth, so the ansatz is not
degenerate; the rank of its outputs survives depolarising and amplitude-damping
noise at rates well above current hardware. The layer could be trained, and it
could be run. It has nothing to contribute here that the classical
encoder around it does not already provide.

\subsection{What a deployment would actually cost}
Because the strategy consumes a ranking, the operative hardware question is how
many shots preserve rank order at the spread of the scores being ranked, and
that spread is set by the trained model, not by the circuit. Across the twenty
trained cells the within-date standard deviation of the circuit output runs
from \num{0.007} to \num{0.51}, with a median of \num{0.10}, because
cross-sectional normalisation and the encoder map a day's securities into a
band of angles whose width the training trajectory decides. Preserving the
ranking to a Spearman correlation of \num{0.99} takes about \num{8192} shots
per security at the median cell, roughly $1.6\times10^{7}$ circuit executions
to score \num{2000} securities for one trading day, and more than
\num{131072} at the narrowest cells; the \num{256} a scan on random inputs
suggests is reached by no cell, and the widest need \num{1024}. Gate noise remains negligible at those budgets
(sampling accounts for essentially all of the degradation), so the cost is set
by measurement statistics alone.

We make no claim about whether a present device sustains that rate, and the
comparison that matters is not with a device specification but with the
classical alternative: the same ranking is produced by a parameter-matched
classical map in microseconds, and by no map at all equally well. The
conclusion is therefore narrower than the one a circuit-level scan would
support. It is not that hardware is no obstacle; it is that the obstacle is
mis-stated when the shot budget is measured on random inputs, or on one
encoder, and that on this task the observable can be concentrated enough to
make ranking expensive even at eight qubits, by an amount that varies a
hundredfold between statistically indistinguishable cells. Whether that
concentration is intrinsic to cross-sectional ranking or an artefact of this
encoder is the question we would put first in
any follow-up.

\subsection{Threats to validity}
\label{sec:discussion:threats}
We list the ways this study could be wrong, in decreasing order of how much
they would change the conclusion.

\paragraph{One feature set, chosen by us.}
All eleven conditions see the same \num{21} technical features. If the quantum
layer has an inductive bias suited to structure these features do not contain
(fundamentals, order-flow, cross-asset signals), the null is a statement
about the inputs rather than the layer. Nothing in the design can rule this
out, and it is the most likely way the result fails to generalise.

\paragraph{Two ansatz families, one of them resolved finely.}
The QCNN (angle embedding, shared two-qubit blocks, pooling, single-wire
measurement) and a re-uploading circuit with a CNOT ring and a readout
averaged over every wire give the same ordering, but only the first is certified equivalent
to its frozen copy inside \num{0.03} Sharpe; the second is bounded at
\num{0.25}, because its all-wire readout lets the parameters reach the
score and the paired standard error grows twelvefold. A family in which
training the circuit helps may exist; neither of these two is it, and the
second shows that a looser bound is what a more expressive circuit buys here.

\paragraph{Low achievable signal.}
The information coefficient available on this task is around \num{0.006}
(Table~\ref{tab:ic}), and the spread across seeds is as large as the
differences between conditions. A layer cannot demonstrate value where the signal
it would improve on is itself indistinguishable from noise. The locked holdout
is the partial answer: on a window where the shared signal does pay, the
ablation is unchanged. It is \num{135} days, which is short by this paper's own
standard. The gradient-boosted baseline is the other partial answer: the model
family that wins this competition extracts an information coefficient of
\num{0.005} from the same features, no more than the networks, so the weak
signal is a property of the inputs rather than of the architectures.

\paragraph{Statistical resolution.}
The design resolves the trained-versus-frozen contrast to a margin of
\num{0.03} Sharpe and the contrasts against the no-circuit control only to
\num{0.50}. The second of those is a weak bound and we do not claim more from
it. Four folds and five seeds give twenty paired cells, which is enough for the
first comparison and not for the others.

\paragraph{Simulation only.}
Every circuit result is from a state-vector simulator. The noise scans are
simulated channels, not device runs, and the shot analysis assumes independent
sampling. No claim here has been checked on hardware.

\paragraph{Survivorship in the auxiliary panels.}
The US universe comes from a current exchange symbol directory and excludes
firms delisted before it. A 2022--2025 extension of the JPX panel was also
assembled from the same source for a delisting audit; no result in this paper
uses it, because it is missing
\num{207} of the \num{2000} universe securities, of which the exchange's
official delisting lists account for all: \num{204} left through takeovers,
management buy-outs or squeeze-outs, one failed the listing-maintenance
standards, one was removed by the exchange and one delisted voluntarily. The
JPX development panel itself has the same character: its \num{2000} securities
are the competition's primary universe, fixed by the organiser at the end of
the training period (Section~\ref{sec:data}). All three universes are common
to every condition and to the earlier result we re-score, so the bias enters
no contrast, but absolute performance levels on any of them are not claimed. The US panel also
required a point-in-time price screen to remove a bankruptcy-reorganisation
splice (Section~\ref{sec:data-us}); we found it by auditing for impossible
returns, and a study that did not audit would not have found it.

\paragraph{Fold count.}
Four walk-forward folds are what the JPX sample supports; a fifth would need
\num{44} more trading days than exist. The US calendar affords five.

\subsection{What would change the conclusion}
A positive result would require a task where the circuit's inductive bias
matches structure in the data that the classical encoder cannot express. The
present data has no such structure to find: over the \num{500} development
days only one condition has a positive gross return, and a transaction cost of
half a basis point removes it; on the \num{135}-day holdout the best
break-even cost is \SI{5.5}{bps}. Claims of quantum
advantage in financial prediction should be accompanied by the control
conditions used here (a frozen circuit and a parameter-matched surrogate at
minimum) and by an evaluation long enough to resolve the effect being
claimed. Where they are, we expect most of them to look like
Table~\ref{tab:headline}.
